\chapter{Podstawy Architektury Aplikacji (Wykład 4)}

\section{Problem Monolitycznego Kodu ("God Composable")}

Na początku przygody z Jetpack Compose naturalnym odruchem jest pisanie całego kodu wewnątrz funkcji \texttt{@Composable}. Skoro UI to tylko funkcja, dlaczego nie umieścić w niej również pobierania danych z sieci czy logiki biznesowej?

Spójrzmy na kod, który działa, ale łamie zasady inżynierii oprogramowania. Nazywamy to antywzorcem \textbf{"God Composable"} - funkcją, która wie i robi wszystko.

\begin{minted}[frame=single,framesep=10pt]{kotlin}
// KROK 1: Kod "naiwny"
@Composable
fun UserProfileScreen() {
    // Stan UI + Logika w jednym miejscu
    var userData by remember { mutableStateOf("Ładowanie...") }
    val scope = rememberCoroutineScope() // Scope powiązany z cyklem życia UI

    Column {
        Button(onClick = {
            // Operacja sieciowa bezpośrednio w UI
            scope.launch {
                try {
                    val user = api.getUser() // Symulacja pobierania (2 sekundy)
                    userData = user.name.uppercase() // Logika biznesowa
                } catch (e: Exception) {
                    userData = "Błąd!"
                }
            }
        }) {
            Text("Pobierz dane")
        }
        Text(userData)
    }
}
\end{minted}

\subsection{Rotacja Ekranu}
Uruchamiamy aplikację, klikamy przycisk, dane się pobierają. Wszystko wygląda świetnie. Do momentu, gdy obrócimy telefon.
Android przy zmianie konfiguracji (np. obrót ekranu) niszczy i tworzy aktywność od nowa. Funkcja \texttt{remember} traci pamięć. Użytkownik, który czekał na dane, nagle znów widzi ekran początkowy \texttt{Ładowanie...}.

\subsection{Próba ratunku: rememberSaveable}
Rozwiązanie tego problemu już znamy z poprzedniego semestru. Spróbujmy więc naprawić nasz kod:

\begin{minted}[frame=single,framesep=10pt]{kotlin}
// rememberSaveable
var userData by rememberSaveable { mutableStateOf("Ładowanie...") }
\end{minted}

Czy to rozwiązuje problem? Tylko pozornie i tylko w trywialnych przypadkach. Gdy spróbujemy zastosować to w większej aplikacji, napotkamy trzy krytyczne problemy, których \texttt{rememberSaveable} nie rozwiąże:

\begin{itemize}
    \item \textbf{Śmierć Korutyny:} \texttt{rememberCoroutineScope} jest ściśle powiązany z widokiem. Gdy obracasz ekran w trakcie pobierania danych (np. wolne WiFi), stary widok ginie, a wraz z nim \textbf{anulowana jest korutyna}. Pobieranie zostaje przerwane w połowie. Nowy ekran powstaje, ale nic nie wie o tym, że poprzedni coś pobierał. Użytkownik klika, czeka i nic nie dostaje.
    
    \item \textbf{Ograniczenia Pamięci (Bundle):} \texttt{rememberSaveable} zapisuje dane w systemowym \texttt{Bundle}. Jest on przeznaczony dla małych danych (tekst, liczby). Jeśli spróbojemy tam zapisać listę 500 obiektów JSON pobranych z API, aplikacja wyrzuci błąd \texttt{TransactionTooLargeException} i się zamknie.
    
    \item \textbf{Testowalność:} Jak przetestować, czy nazwisko jest poprawnie zamieniane na wielkie litery? W obecnym kodzie jest to niemożliwe bez uruchamiania emulatora, bo logika jest \textit{zabetonowana} wewnątrz przycisku.
\end{itemize}

\subsection{Separacja odpowiedzialności}
Dochodzimy do wniosku, że funkcja \texttt{Composable} nie jest odpowiednim miejscem na trzymanie danych ani wykonywanie operacji. Potrzebujemy miejsca, które:
\begin{enumerate}
    \item Przeżyje rotację ekranu (nie jak \texttt{rememberSaveable}).
    \item Pozwoli korutynom dokończyć pracę, nawet gdy widok jest niszczony.
    \item Nie ma limitu wielkości danych (nie jak \texttt{Bundle}).
\end{enumerate}

\subsection{Metafora Budowlana: Buda vs Wieżowiec}
Aby zrozumieć potrzebę architektury, posłużmy się analogią budowlaną:
\begin{itemize}
    \item \textbf{Budowanie budy dla psa (Brak architektury):} Możesz to zrobić sam. Jeśli wbijesz gwóźdź w złym miejscu, łatwo to poprawić. To są małe, proste aplikacje.
    \item \textbf{Budowanie wieżowca (Z architekturą):} Tutaj potrzebujesz planu. Elektryk (Data Layer) nie maluje ścian, a malarz (UI Layer) nie kładzie instalacji gazowej. W dużym projekcie, jeśli Composable zajmuje się logiką API, to tak jakby malarz próbował naprawiać windę.
\end{itemize}

Potrzebujemy więc "Kierownika Budowy", który ma plany w biurze (bezpiecznym od remontów) i zarządza pracami. Tym kierownikiem będzie komponent, który za chwilę poznamy.

\section{Ewolucja Wzorców Architektonicznych}

Skoro wiemy już, że wrzucanie wszystkiego do jednego worka ("God Composable") jest nie najlepszym rozwiązaniem, musimy zastanowić się, jak podzielić aplikację. W inżynierii oprogramowania ten problem nie jest nowy. Przez dekady wykształciło się wiele różnych podejść do tego problemu. Jednym z nich jest rodzina wzorców MVx, które różnią się sposobem, w jaki warstwa danych komunikuje się z warstwą wizualną.

Aby zrozumieć różnice między nimi, posłużymy się analogiami z życia codziennego.

\subsection{MVC (Model-View-Controller) - Sklep z ladą}


[Image of MVC architecture diagram]

Jest to najstarsze podejście, które można porównać do wizyty w tradycyjnym sklepie, gdzie towar podaje sprzedawca.

\begin{itemize}
    \item \textbf{Ty (View):} Stoisz przed ladą. Widzisz towar, ale jesteś pasywny - nie możesz go sam wziąć.
    \item \textbf{Sprzedawca (Controller):} To on rządzi procesem. Mówisz mu: \textit{Poproszę chleb} (Interakcja).
    \item \textbf{Magazyn (Model):} Sprzedawca idzie na zaplecze, sprawdza stan magazynowy i bierze towar.
\end{itemize}

\textbf{Wniosek:} W tym układzie Sprzedawca (Controller) decyduje o wszystkim. To on musi wiedzieć, jak wygląda magazyn i to on decyduje, co pokazać klientowi. Musi też być w stanie wskazać konkretnego klienta, któremu musi wydać towar.

\subsection{MVP (Model-View-Presenter) - Restauracja}

Druga odmiana wzorca z rodziny MVx.

\begin{itemize}
    \item \textbf{Klient przy stoliku (View):} Siedzisz i czekasz. Nie wiesz i nieinteresuje Cię, co się dzieje w kuchni.
    \item \textbf{Kelner (Presenter):} Przyjmuje zamówienie i idzie do kuchni (Model).
    \item \textbf{Kluczowa cecha:} Kelner wraca z kuchni i \textbf{własnoręcznie stawia talerz na Twoim stole}.
\end{itemize}

W kodzie wygląda to tak, że Presenter (Kelner) posiada referencję do Widoku (Klienta). Wywołuje na nim konkretną metodę, np. \texttt{view.showDinner()}.
\textbf{Problem:} Jest to sztywne połączenie 1:1. Kelner musi znać konkretny stolik. Jeśli Klient wyjdzie do toalety (Rotacja Ekranu/Zniszczenie Widoku), a Kelner wróci z talerzem i spróbuje go postawić na pustym miejscu, dojdzie do błędu (\texttt{NullPointerException}).

\subsection{MVVM (Model-View-ViewModel) - Automat z napojami}

To podejście jest standardem we współczesnym Androidzie i idealnie współgra z Jetpack Compose.

\begin{itemize}
    \item \textbf{Użytkownik (View):} Podchodzi do maszyny i naciska przycisk (Wysyła Event).
    \item \textbf{Automat (ViewModel):} Odbiera sygnał, mieli kawę, pobiera wodę (komunikuje się z Repozytorium/Modelem).
    \item \textbf{Wynik (State):} Automat wystawia kubek z kawą do szuflady odbiorczej (\textit{Emisja stanu}).
\end{itemize}

\textbf{Kluczowa różnica - Reaktywność:} Automatu nie obchodzi, kto ten kubek weźmie.
\begin{itemize}
    \item Automat tylko \textbf{wystawia stan}.
    \item Jeśli nikt nie stoi przed maszyną - kubek (dane) po prostu czeka w szufladzie.
    \item Jeśli użytkownik odejdzie (obróci ekran) i przyjdzie nowy (\texttt{Activity} zostanie stworzone na nowo) - kubek z kawą nadal tam jest. Nowy widok po prostu zagląda do szuflady i widzi gotowy napój.
\end{itemize}

W Compose tą \textit{szufladą} są strumienie danych, takie jak \texttt{StateFlow}. ViewModel aktualizuje \texttt{StateFlow}, a Widok (Composable) jedynie nasłuchuje zmian. To rozwiązuje nasz problem z rotacją ekranu - ViewModel trzyma \textit{napój}, niezależnie od tego, co dzieje się z \textit{klientem}.

\subsection{MVI (Model-View-Intent) - Kiosk samoobsługowy}

Jest to ewolucja wzorca MVVM, która kładzie rygorystyczny nacisk na jednokierunkowy przepływ danych (Unidirectional Data Flow).

\begin{itemize}
    \item \textbf{View (Ekran kiosku):} Wyświetla aktualny stan zamówienia. Jest to \textit{tylko odczyt}.
    \item \textbf{Intent (Zamiar/Akcja):} Gdy klikasz \textit{Dodaj frytki}, nie zmieniasz bezpośrednio zamówienia. Wysyłasz do systemu \textit{zamiar} (Intent) o treści: "Użytkownik chce dodać frytki".
    \item \textbf{Model (System):} System bierze Twój \textit{aktualny paragon}, bierze Twój \textit{zamiar}, przetwarza je i wypluwa \textbf{zupełnie nowy, zaktualizowany paragon} (State).
\end{itemize}

\textbf{Kluczowa cecha - Cykliczność i Niezmienność:}
W przeciwieństwie do MVVM, gdzie ViewModel może mieć wiele różnych strumieni danych (osobno imię, osobno lista, osobno błędy), w MVI dążymy do posiadania \textbf{jednego obiektu stanu} (np. \texttt{OrderUiState}).
W naszej metaforze: Nie możesz wziąć długopisu i dopisać frytek do wydrukowanego paragonu. Musisz wysłać żądanie, a system wydrukuje Ci nowy, zaktualizowany paragon. To zapewnia przewidywalność aplikacji - zawsze wiemy, co doprowadziło do obecnego stanu ekranu.

W Jetpack Compose naturalnie dążymy do MVI, używając pojedynczego \texttt{StateFlow<UiState>} w ViewModelu.